\section*{Problem 6: Linear Programming on Multidimensional Knapsack}
In the Multidimensional Knapsack problem we are given a d-dimensional
knapsack, for some $d \geq 2$, having capacity $b_i$ in the $i$th
dimension, for $1 \leq i \leq d$. Also, there is a set of items $I =
\{1, \ldots, n\}$. The size of item $j$ in dimension $i$ is $s_{ij} >
0$, $1 \leq j \leq n$, $1 \leq i \leq d$. The profit from packing
item $j$ is $p_j > 0$. A solution is a subset of items $A \subseteq
I$, and the profit of $A$ is $P(A) = \sum_{j \in A} p_j$.
A solution $A$ is feasible if $\sum_{j \in A} s_{ij} \leq b_i$ for
all $1 \leq i \leq d$.
We assume below that $d$ is a fixed constant. Consider the following
algorithm for $d$-dimensional Knapsack. Let $\epsilon > 0$ and $h =
\min(n, \lceil d(1+\epsilon)/\epsilon \rceil)$. Given a selected
subset of items $S \subseteq I$, the objective of the $S$-problem is
to find a feasible packing of items in the knapsack which maximizes
the profit, subject to the following constraints:
(i) All items in $S$ are packed.
(ii) Among the items in $I \setminus S$ only items $j$ with profit
$p_j \leq \min_{\ell \in S} p_\ell$ can be added to the solution.
Denote this set of items by $\text{Good}(S)$.
Given the selected subset of items $S$, we denote by LP (S) a linear
programming relaxation of the $S$-problem. Let $x_j$ be an indicator
for the selection of item $j \in I$ for the solution. In the linear
programming relaxation $LP(S)$, $0 \leq x_j \leq 1$. Let $x^B =
(x^B_1, \ldots, x^B_n)$ be a basic solution for $LP(S)$.

\textbf{The Algorithm}: For any $S \subseteq I$ such that $|S| \leq
h$, solve optimally $LP(S)$ and round down $x^B$. Choose the solution
which maximizes the profit.
(a) (2 pt.) Formulate the linear program $LP(S)$.
(b) (5 pt.) Denote by $x^I_j$ the integer variable obtained after
rounding the basic variable $x^B_j$. Let $F$ be the subset of items
$j$ for which $x^B_j$ is fractional, i.e., $x^I_j < x^B_j$. Show that
$|F| \leq d$.

\subsection*{Formulation of LP(S)}
The linear programming relaxation $LP(S)$ can be formulated as follows:

\textbf{Maximize:}
\[ \sum_{j \in I} p_j x_j \]

\textbf{Subject to:}
\begin{align*}
    \sum_{j \in I} s_{ij} x_j &\leq b_i && \forall i \in \{1, \ldots, d\} \\
    x_j &= 1 && \forall j \in S \\
    x_j &= 0 && \forall j \in I \setminus (S \cup \text{Good}(S)) \\
    0 \leq x_j &\leq 1 && \forall j \in \text{Good}(S)
\end{align*}


\subsection*{Bounding the number of fractional variables}
Let $x^B$ be a basic feasible solution to $LP(S)$. There are $n$ variables in the linear program. In a basic feasible solution, there must be exactly $n$ linearly independent tight (active) constraints.

From the formulation of $LP(S)$, we have the following constraints that can be tight:
\begin{enumerate}
    \item $x_j = 1$ for $j \in S$ (exactly $|S|$ constraints)
    \item $x_j = 0$ for $j \in I \setminus (S \cup \text{Good}(S))$ (exactly $n - |S| - |\text{Good}(S)|$ constraints)
    \item $x_j = 0$ or $x_j = 1$ for $j \in \text{Good}(S)$
    \item $\sum_{j \in I} s_{ij} x_j \leq b_i$ for $i \in \{1, \ldots, d\}$ (at most $d$ constraints)
\end{enumerate}

Let $k$ be the number of variables $x_j$ for $j \in \text{Good}(S)$ that are tight at their bounds (i.e., $x_j = 0$ or $x_j = 1$).
Since a basic feasible solution is determined by $n$ linearly independent tight constraints, we must have:
\[ |S| + (n - |S| - |\text{Good}(S)|) + k + (\text{number of tight capacity constraints}) \geq n \]
This simplifies to:
\[ n - |\text{Good}(S)| + k + (\text{number of tight capacity constraints}) \geq n \]
\[ k \geq |\text{Good}(S)| - (\text{number of tight capacity constraints}) \]

Since there are only $d$ capacity constraints in total, the number of tight capacity constraints is at most $d$. Thus:
\[ k \geq |\text{Good}(S)| - d \]

This means that out of the $|\text{Good}(S)|$ variables in $\text{Good}(S)$, at least $|\text{Good}(S)| - d$ of them are fixed to $0$ or $1$. The remaining variables are the only ones that can be strictly fractional ($0 < x_j < 1$).
Therefore, the number of fractional variables is at most:
\[ |F| = |\text{Good}(S)| - k \leq |\text{Good}(S)| - (|\text{Good}(S)| - d) = d \]
Therefore, $|F| \leq d$.
